% Four candidate explanations for the failure of probe relocation, each with its test and result.
% Written for the paper from results/c2_analysis/ and results/supplementary/ (mlp_probe.csv, globalisation.csv)
% as summarised in paper/evidence_dossier.md; no number is computed here.
\footnotesize
\setlength{\tabcolsep}{4pt}
\renewcommand{\arraystretch}{1.08}
\begin{tabular}{>{\raggedright\arraybackslash}p{3.1cm}>{\raggedright\arraybackslash}p{4.9cm}>{\raggedright\arraybackslash}p{5.1cm}}
\toprule
Explanation & Test (layers 9--12 unless stated) & Result \\
\midrule
The probe sends score to distant, wrong patches & Share of the probe's error mass more than three patches from the true patch & Small: 2--29\% across models and layers. At layer 12, 45\%, 66\% and 81\% of the probability lies within the $3{\times}3$ neighbourhood of the true patch \\
The probe has a fixed position preference & Peak of the probe's mean prediction relative to uniform; its correlation with the mean map & Close to uniform (peak 1.1--1.3 times the uniform value); at most weakly related to the mean map \\
The linear probe is too weak & Two-layer MLP probe in place of the linear one & Accuracy at layer 12 rises by only 1.8, 3.8 and 6.7 points; relocation with the MLP probe is still lower than received attention in all six cells \\
Global tokens are harder to place & Probe accuracy and entropy against a token's similarity to the image's mean token (layers 10--12) & No consistent relation; accuracy falls with globality only in DeiT-B \\
\bottomrule
\end{tabular}
